\documentclass[11pt]{amsart}

\usepackage[T1]{fontenc}
\usepackage[utf8]{inputenc}
\usepackage{microtype}
\usepackage{mathtools}
\usepackage{amssymb}
\usepackage{comment}
\usepackage{enumitem}
% Shared review markup and formal links; the human draft hides declaration names.
% Review annotations, following the red-add/blue-remove convention used in
% the Varenna 2016 papers. These commands keep the original wording visible.
% Accept an insertion by removing its AIaddition environment or AIadd wrapper.
% Accept a replacement by deleting AIremove{old} and unwrapping AIadd{new}.
% Resolve a comment by deleting its complete AIcomment{ID}{text} command.
% There is deliberately no global switch that hides unresolved comments.
\usepackage{xcolor}
\usepackage{xurl}
\usepackage[normalem]{ulem}
% Give long inline review markup a little extra line-breaking room without
% changing the author text or the mathematical content.
\setlength{\emergencystretch}{.75em}
% AI compile repair: the retained draft uses the undefined typo \cots.
% EDIT-02 beside that occurrence makes this repair visible in the PDF.
\providecommand{\cots}{\cdots}
% Permit a line break after commas in long inline coordinate lists.
\begingroup\lccode`~=`,
\lowercase{\endgroup\def~}{\mathchar"613B\allowbreak}
\mathcode`,="8000
\definecolor{AIadded}{RGB}{160,25,25}
\definecolor{AIreview}{RGB}{25,70,145}
\newcommand{\AIadd}[1]{{\color{AIadded}#1}}
\newcommand{\AIremove}[1]{{\color{AIreview}%
  \ifmmode\text{\sout{\ensuremath{#1}}}\else\sout{#1}\fi}}
\newcommand{\AIreplace}[2]{\AIremove{#1}\,\AIadd{#2}}
\newenvironment{AIaddition}[1]
  {\par\begingroup\color{AIadded}\noindent
   \textbf{AI addition: #1.}\par\smallskip\ignorespaces}
  {\par\endgroup}
\newcommand{\AIcomment}[2]{%
  \par\smallskip\noindent
  \fcolorbox{AIreview}{AIreview!4}{%
    \parbox{\dimexpr\linewidth-2\fboxsep-2\fboxrule\relax}{%
      \normalfont\small\color{AIreview}\textbf{AI comment #1.} #2}}%
  \par\smallskip}
\newcommand{\AIreviewlegend}{%
  \AIcomment{Review guide}{Red text proposes additions or corrections;
  blue struck text preserves the original wording.
  \textbf{Johanna:} review the proposals as mathematics and prose, and resolve
  their boxes after accepting or revising them. No Lean installation, code
  editing, or proof replay is required for her review.
  \textbf{Max:} review the whole paper--Lean correspondence, including
  every statement, its hypotheses, and the definitions entering the challenge. Comparison
  notes explain proof differences beside the relevant passages; the technical
  supplement gives mathematical accounts and optional declaration links. The
  side-by-side review interface links statements to their checked sources.}}

% Lean cross-references for the paper--Lean audit.
% Each formalized result is followed by \formalref{...}; every \leandecl or
% \leanlib names one declaration by file, line and fully qualified name at
% the pinned revisions below. The review generator in itpplasma/stafford38-formal
% (tools/paper_lean_audit) checks every such reference against the Lean
% source, so a stale file, line or name fails its build.
% Accepting: keep this file and the \formalref lines, or delete both.
\newcommand{\leanrepo}{https://github.com/itpplasma/stafford38-formal/blob/591d4e11e3971577be169e428ac73c44a4ac8a32}
\newcommand{\leanlibrepo}{https://github.com/itpplasma/algebraic-analysis/blob/4aae47967f6ba02ffe2f639ab06564c9a9d1ecc8}
% \leandecl{file}{line}{Fully.Qualified.Name}: declaration in stafford38-formal.
\newcommand{\leandecl}[3]{\href{\leanrepo/#1\#L#2}{\nolinkurl{#3}}}
% \leanlib{file}{line}{Fully.Qualified.Name}: declaration in AlgebraicAnalysis.
\newcommand{\leanlib}[3]{\href{\leanlibrepo/#1\#L#2}{\nolinkurl{#3}}}
% The human-readable paper displays only substantive mathematical qualifications.
% Full file paths and declaration names remain in source and in the supplement.
\newif\ifshowleandetails
\showleandetailsfalse
% \formalref[relation to the printed statement]{declarations}
\newcommand{\formalref}[2][]{%
  \ifshowleandetails
    \par\smallskip\noindent
    {\small\textsf{Lean}\if\relax\detokenize{#1}\relax\else\ (#1)\fi: #2}%
    \par\smallskip
  \else
    \if\relax\detokenize{#1}\relax\else
      \par\smallskip\noindent{\small\textsf{Formal comparison:} #1}%
      \par\smallskip
    \fi
  \fi}

\usepackage{xr-hyper}
\usepackage[colorlinks=true,citecolor=blue,linkcolor=blue,urlcolor=blue]{hyperref}
\numberwithin{equation}{section}

\newtheorem{theorem}{Theorem}[section]
\newtheorem{proposition}[theorem]{Proposition}
\newtheorem{lemma}[theorem]{Lemma}
\newtheorem{corollary}[theorem]{Corollary}
\theoremstyle{plain}
\newtheorem{definition}[theorem]{Definition}
\newtheorem{convention}[theorem]{Convention}
\theoremstyle{remark}
\newtheorem{remark}[theorem]{Remark}



\title{Stafford 3.8: technical formal-proof comparisons}
\author{Christopher Albert et al.}
\address{Graz University of Technology}
\date{}
\newcommand{\Char}{\operatorname{Char}}
\newcommand{\Weyl}{A_n}
\newcommand{\ord}{\operatorname{ord}}
\newcommand{\gr}{\operatorname{gr}}
\newcommand{\Ann}{\operatorname{Ann}}
\newcommand{\Dir}{\operatorname{Dir}}
\newcommand{\Sp}{\operatorname{Sp}}
\newcommand{\Spec}{\operatorname{Spec}}
\newcommand{\Supp}{\operatorname{Supp}}
\newcommand{\Sing}{\operatorname{Sing}}
\newcommand{\sm}{\mathrm{sm}}
\newcommand{\bbA}{\mathbb{A}}
\newcommand{\bbP}{\mathbb{P}}
\newcommand{\bbN}{\mathbb{N}}
\newcommand{\bbZ}{\mathbb{Z}}
\newcommand{\bbQ}{\mathbb{Q}}
\newcommand{\Kbar}{\overline{k}}
\newcommand{\dd}{\mathrm{d}}
\DeclareMathOperator{\rad}{rad}


\externaldocument[paper:]{human_readable_main}[https://itpplasma.github.io/stafford38-formal/human_readable_main.pdf]
% Declaration paths are hidden in the reader copy. Uncomment to display them.
% \showleandetailstrue
\begin{document}
\maketitle
% Readable mathematical comparisons supplement the concise COMP annotations
% in the human-readable manuscript.
\AIcomment{Technical review}{\textbf{Information --- readable proof accounts.} This supplement explains the formal
arguments as mathematics. These accounts are AI additions for author review. References
point to the human-readable paper; this supplement is not required reading
for editing that paper.}

\par\bigskip\noindent\pdfbookmark[1]{Comparison COMP-01}{bookmark-COMP-01}\hypertarget{comparison-COMP-01}{\textbf{Comparison COMP-01}}\label{comparison-COMP-01}\par\smallskip
\AIcomment{COMP-01}{\textbf{Information --- no paper edit requested.} Lean follows the printed localization proof: extend the bracket by the quotient rule, use minimality to identify the localized radical with the maximal ideal, and contract. Its bracket has the opposite sign; involutivity is unchanged. \href{https://itpplasma.github.io/stafford38-formal/lean_proof_details.pdf\#nameddest=comparison-COMP-01}{Mathematical proof (COMP-01, p.~1).}}
\begin{AIaddition}{Proof as formalized in Lean}
Lean proves the lemma in the corrected form of ALG-01. Let $J$ be any ideal
of $R=k[x,\xi]$ with $\{\sqrt J,\sqrt J\}\subseteq\sqrt J$, and let
$\mathfrak p$ be minimal over $J$. Then
$\{\mathfrak p,\mathfrak p\}\subseteq\mathfrak p$. The official proof
follows the paper's localization argument. The partial derivatives extend
to $R_{\mathfrak p}$ by the quotient rule
\[
\partial_i(a/s)=(s\partial_i a-a\partial_i s)/s^2.
\]
Their Poisson bracket is a biderivation and preserves the extension of
$\sqrt J$. Minimality gives
$(\sqrt J)R_{\mathfrak p}=\mathfrak pR_{\mathfrak p}$, so this maximal
ideal is involutive. Contracting it to $R$ proves the assertion. Lean's
bracket has the opposite sign from the displayed paper convention;
involutivity is unchanged.

The official Gabber interface consumes the localized proof above to pass
from radical involutivity to minimal-prime involutivity.
\end{AIaddition}

\par\bigskip\noindent\pdfbookmark[1]{Comparison COMP-02}{bookmark-COMP-02}\hypertarget{comparison-COMP-02}{\textbf{Comparison COMP-02}}\label{comparison-COMP-02}\par\smallskip
\AIcomment{COMP-02}{\textbf{Information --- no paper edit requested.} Lean checks the printed induction in an arbitrary symplectic space, including completion with the prescribed first vector. The official normalization uses the resulting basis-to-basis map and its matrix. \href{https://itpplasma.github.io/stafford38-formal/lean_proof_details.pdf\#nameddest=comparison-COMP-02}{Mathematical proof (COMP-02, p.~1).}}
\begin{AIaddition}{Proof as formalized in Lean}
Lean follows the printed basis induction in an arbitrary finite-dimensional
space with an alternating nondegenerate form. For a prescribed nonzero $v$,
choose $w$ with $\omega(v,w)=1$. The plane $U=\operatorname{span}(v,w)$
is nondegenerate; Lean proves $\mathbb V=U\oplus U^\perp$, nondegeneracy
on $U^\perp$, and its dimension $2n-2$. Induction gives a symplectic basis
of the complement and hence a basis beginning with $v,w$.

Completing each of two nonzero vectors gives bases with the same pairwise
form values. The basis-to-basis map therefore preserves the form and sends
one prescribed vector to the other. The official normalization caller
converts this map into its matrix and checks the standard-form identity.
The abstract basis pairing and the coordinate matrix use opposite Gram
sign conventions; the conversion proves the matrix identity explicitly.
No change to the paper's basis proof is needed.
\end{AIaddition}

\par\bigskip\noindent\pdfbookmark[1]{Comparison COMP-03}{bookmark-COMP-03}\hypertarget{comparison-COMP-03}{\textbf{Comparison COMP-03}}\label{comparison-COMP-03}\par\smallskip
\AIcomment{COMP-03}{\textbf{Information --- proof comparison; no edit requested.} The formal normalization uses an explicit symplectic matrix and its inverse. Its symbol substitution is $s(MX)$; with the convention used here this is $s\circ g^{\mathsf T}$. Monicity means that the $p^N$ PBW coefficient is $1$ and the Bernstein degree is at most $N$, which gives the coefficient bounds above. \href{https://itpplasma.github.io/stafford38-formal/lean_proof_details.pdf\#nameddest=comparison-COMP-03}{Mathematical proof (COMP-03, p.~2).}}
\begin{AIaddition}{Proof as formalized in Lean}
Lean writes the change of generators with a matrix $M=g^{\mathsf T}$:
$\Phi_M(z_i)=\sum_jM_{ij}z_j$, where
$z=(x_1,\ldots,x_n,\partial_1,\ldots,\partial_n)$. Put $J_{ij}:=[z_i,z_j]$. Then
$[\Phi_Mz_i,\Phi_Mz_j]=(MJM^{\mathsf T})_{ij}$, so $\Phi_M$ is an algebra
endomorphism whenever $MJM^{\mathsf T}=J$. If $M'$ also satisfies this and
$MM'=M'M=1$, then $\Phi_{M'}$ is the inverse of $\Phi_M$. For
$a\in B_N$ one has $\Phi_M(a)\in B_N$ and
$\sigma^B_N(\Phi_Ma)(X)=\sigma^B_N(a)(MX)$, where
$(MX)_i=\sum_jM_{ij}X_j$.

Let $N$ be the Bernstein degree of $d$, i.e.\ the top total degree of its PBW
normal form. Let $s=\sigma^B_N(d)\neq0$; it is homogeneous of degree $N$. If
$N=0$, $d$ is a nonzero scalar. Let $N\ge1$ and let $t$ be the index of
$\partial_1$. The argument has five steps.

(i) Since $k$ is infinite, there is $v\in k^{2n}$ with $s(v)\neq0$. Here
$v\neq0$, because $s(0)=0$ for $N\ge1$.

(ii) Lemma~\ref{paper:lem:sympcompl} with $e=e_t$ gives a symplectic $M$ whose
$t$-th column is $v$, i.e.\ $Me_t=v$. Its inverse $M'$ is symplectic too.

(iii) The symbol $s_M(X):=s(MX)$ of $\Phi_M(d)$ is homogeneous of degree
$N$. On the $\xi_1$-axis it restricts to $c\,\xi_1^N$, where $c$ is its
coefficient of $\xi_1^N$. Hence $c=s_M(e_t)=s(Me_t)=s(v)\neq0$.

(iv) Consequently $\tilde d:=c^{-1}\Phi_M(d)$ lies in $B_N$, and the
coefficient of $\xi_1^N$ in $\sigma^B_N(\tilde d)$ is $1$. This records the
chosen symplectic change of coordinates together with its normalized leading
coefficient.

(v) $\sigma^B_N(\tilde d)$ is the weight-$N$ part of the PBW normal form, so
this coefficient is the coefficient of the PBW word $\partial_1^N$ in
$\tilde d$. The resulting normal form has these properties: every PBW word
$x'^{\alpha}\partial'^{\beta}x_1^{a}p^{b}$ of $\tilde d$ has
$|\alpha|+|\beta|+a+b\le N$, and $p^N$ has coefficient $1$. The words are
ordered transverse pairs first, then $x_1^ap^b$, with $p$ rightmost. Grouping
by $b$ gives
$\tilde d=p^N+\sum_{j<N}c_jp^j$ with $c_j\in C$ and $\deg_Bc_j\le N-j$, with
the coefficients on the \emph{left}. Lean records this as monicity of the
outer Ore polynomial.

Lean additionally checks the paper's right-coefficient form
$\tilde d=p^N+\sum_{j<N}p^ja_j$. Moving coefficients across $p$ uses the
coefficient derivation, which lowers Bernstein degree. The right PBW basis
therefore gives $\deg_Ba_j\le N-j$, top coefficient $a_N=1$, and zero
coefficients above $N$. Its finite-support reconstruction proves the
stated equality. The paper can keep its original right-coefficient convention.
\end{AIaddition}

\par\bigskip\noindent\pdfbookmark[1]{Comparison COMP-04}{bookmark-COMP-04}\hypertarget{comparison-COMP-04}{\textbf{Comparison COMP-04}}\label{comparison-COMP-04}\par\smallskip
\AIcomment{COMP-04}{\textbf{Information --- no paper edit requested.} Lean checks the Euler weights of PBW monomials and identifies the nonnegative Euler part with the transverse--Euler subring used in the residue argument. The corrected product identities and one-sided normality give the same surjectivity conclusion. \href{https://itpplasma.github.io/stafford38-formal/lean_proof_details.pdf\#nameddest=comparison-COMP-04}{Mathematical proof (COMP-04, p.~2).}}
\begin{AIaddition}{Proof as formalized in Lean}
Lean checks the Euler weights of PBW monomials directly:
$[E,bx_1^ap^j]=(a-j)bx_1^ap^j$ for transverse $b$.
It proves that the subring generated by the nonnegative-weight PBW monomials
is exactly the subring $R\subseteq A_n$ generated by
$T'=A_{n-1}(k)$, $x_1$ and $E=x_1p$. This identifies the paper's
nonnegative Euler part with the subring used by the residue proof.
Internally, Lean's Ore normal form puts the coefficients on the
left: every element is uniquely $\sum_jc_jp^j$ with $c_j\in C=T'[x_1]$, that
is, a sum of words $b\,x_1^ap^j$ with $b\in T'$. The proof has five steps.

(a) $R$ contains $\bbQ[E]$. It is \emph{one-sidedly normal}: for
$r\in R$ there is $r'\in R$ with $x_1r=r'x_1$. The set of such $r$ is a
subring, and it contains the generators, because $x_1b=bx_1$,
$x_1x_1=x_1x_1$ and $x_1E=(E-1)x_1$.

(b) The polynomials $\prod_{i=0}^{N-1}(X-i)$ and $\prod_{i=1}^{N}(X+i)$ have
the disjoint root sets $\{0,\ldots,N-1\}$ and $\{-1,\ldots,-N\}$. They are
therefore coprime in $\bbQ[X]$. With Lemma~\ref{paper:lem:euler} (corrected
form) this gives $u,v\in\bbQ[X]$ with
$p^Nx_1^N\,u(E)+x_1^Np^N\,v(E)=1$.

(c) The normal form of Proposition~\ref{paper:prop:monic} gives
$\tilde d=p^N+\sigma$, where $\sigma$ is a sum of words $b\,x_1^ap^j$ with
$j<N$. Then
\[
\tilde d\,\bigl(x_1^Nu(E)\bigr)+\bigl(x_1^N\tilde d\bigr)v(E)
=1+\sigma x_1^Nu(E)+x_1^N\sigma v(E),
\]
and the left side lies in $I$.

(d) Each word $w=b\,x_1^ap^j$ with $j<N$ has a right factor $x_1$ on both
sides, with cofactors in $R$. By Lemma~\ref{paper:lem:euler} and (a),
\begin{align*}
w\,x_1^N&=b\,x_1^a(p^jx_1^j)\,x_1^{N-j}
=\Bigl[b\,x_1^a\prod_{i=1}^{j}(E+i)\,x_1^{N-j-1}\Bigr]x_1,\\
x_1^Nw&=b\,x_1^{N+a-j-1}\,x_1\,(x_1^jp^j)
=\Bigl[b\,x_1^{N+a-j-1}\prod_{i=0}^{j-1}(E-1-i)\Bigr]x_1 .
\end{align*}
By additivity, $\sigma x_1^N=U'x_1$ and $x_1^N\sigma=V'x_1$ with
$U',V'\in R$.

(e) Since $x_1u(E)=u(E-1)x_1$, the residue equals $Ux_1$ with
$U=U'u(E-1)+V'v(E-1)\in R$. Hence $1+Ux_1\in I$.
\end{AIaddition}

\par\bigskip\noindent\pdfbookmark[1]{Comparison COMP-05}{bookmark-COMP-05}\hypertarget{comparison-COMP-05}{\textbf{Comparison COMP-05}}\label{comparison-COMP-05}\par\smallskip
\AIcomment{COMP-05}{\textbf{Information --- proof comparison; no edit requested.} The formal proof works element by element with the transverse--Euler subring. The residue identity gives $1+Ux_1\in I$ for a suitable $U$ in that subring; one-sided normality and the monic relation then give surjectivity of right multiplication by $x_1$. This avoids treating a quotient by a vector subspace as an algebra quotient. \href{https://itpplasma.github.io/stafford38-formal/lean_proof_details.pdf\#nameddest=comparison-COMP-05}{Mathematical proof (COMP-05, p.~3).}}
\begin{AIaddition}{Proof as formalized in Lean}
Lean argues element by element instead of through the submodule
$qA_{\ge0}$. Let $R$, $U$ be as in Lemma~\ref{paper:lem:euler-residue} (Lean
form), and write $a\equiv b$ for $a-b\in I$. Put
\[
\mathcal P:=\{a\in A_n:\ \text{for every }r\in R\text{ there is }r'\in R
\text{ with } ra\equiv r'\}.
\]
$\mathcal P$ is a subring containing $R$. It is closed under products:
if $ra\equiv s$ and $sb\equiv t$, then $rab\equiv sb\equiv t$, because $I$ is
a right ideal. It also contains $p$. Indeed, for $r\in R$ pick $r'\in R$ with
$x_1r=r'x_1$. From $(1+Ux_1)rp\in I$ we get $rp\equiv-Ux_1rp=-Ur'E\in R$.

The Ore normal form shows that $T'$, $x_1$ and $p$ generate $A_n$, and
$T',x_1\in R$. Hence $\mathcal P=A_n$. Given $a\in A_n$, take $r=1$: then
$a\equiv r_a$ for some $r_a\in R$. Choosing $r'$ with $x_1r_a=r'x_1$ gives
$a\equiv r_a\equiv-Ux_1r_a=(-Ur')x_1$. So the class of $a$ is the class of
$-Ur'$ multiplied on the right by $x_1$.

Lean proves this in the iterated Ore model of $A_n$ and transports it along
the isomorphism with the presented Weyl algebra.
\end{AIaddition}

\par\bigskip\noindent\pdfbookmark[1]{Comparison COMP-06}{bookmark-COMP-06}\hypertarget{comparison-COMP-06}{\textbf{Comparison COMP-06}}\label{comparison-COMP-06}\par\smallskip
\begin{AIaddition}{Proof as formalized in Lean}
Lean never localizes $Q$; only the pages carry a $T$-module structure. Put
$G^p:=F_{-p}Q$ for $p\le0$ and $G^p:=0$ for $p>0$. This is a decreasing,
exhaustive filtration with $G^1=0$. The complex is the $k$-linear map
$f=R_{x_1}\colon Q\to Q$, with $f(G^p)\subseteq G^p$. For $r\ge0$ define
\[
 A_r^p=\frac{G^p\cap f^{-1}(G^{p+r})}{G^{p+1}\cap f^{-1}(G^{p+r})},
 \qquad B_r^p=\frac{G^p}{(G^p\cap f(G^{p-r+1}))+G^{p+1}},
\]
and $A_r=\bigoplus_{p\in\mathbb Z} A_r^p$, $B_r=\bigoplus_{p\in\mathbb Z} B_r^p$.
Applying $f$ induces $d_r\colon A_r^p\to B_r^{p+r}$: changing a representative
by the source denominator changes its image by the target denominator.
Thus $A_0=B_0=M$ and $d_0$ is multiplication by $x_1$. The proof has four steps.

\emph{Step 1: $T$ acts on the pages.} Let $a$ be a transverse generator
$x_j$ or $\partial_j$ ($j\ge2$), of order $e_a\in\{0,1\}$. Then $R_a$ is
$k$-linear, $R_a(G^p)\subseteq G^{p-e_a}$, and $R_a$ commutes with $f$
because $ax_1=x_1a$. It therefore induces maps $A_r^p\to A_r^{p-e_a}$ and
$B_r^p\to B_r^{p-e_a}$ that commute with $d_r$. For two generators,
$[R_a,R_b]=\pm R_{[a,b]}$ with $[a,b]\in F_{e_a+e_b-1}$, so it maps $G^p$ into
$G^{p-e_a-e_b+1}$. An $f$-commuting operator with this extra drop induces
$0$ on every page, because its image lies in the denominator. The induced
operators therefore commute. Letting $x_j,\xi_j\in T$ act by $R_{x_j}$ and
$R_{\partial_j}$ makes every $A_r$, $B_r$ a $T$-module and every $d_r$
$T$-linear. For $r=0$ this is the $T$-structure of $M=\gr Q$, so
$A_1\cong V$ and $B_1\cong U$ as $T$-modules. In addition,
$A_{r+1}\cong\ker d_r$ and $B_{r+1}\cong\operatorname{coker}d_r$ as
$T$-modules. Indeed, if $d_r([z])=0$, write $f(z)=f(w)+v$ with
$w\in G^{p+1}$ and $v\in G^{p+r+1}$. Then $z-w$ represents the same source
class and lies in $G^p\cap f^{-1}G^{p+r+1}$; its kernel of passage back
is exactly the denominator of $A_{r+1}^p$. For the target, quotienting
$B_r^{p+r}$ by $d_r(A_r^p)$ enlarges the boundary denominator from
$G^{p+r}\cap f(G^{p+1})$ to $G^{p+r}\cap f(G^p)$, giving $B_{r+1}^{p+r}$.
Consequently $0\to A_{r+1}\to A_r\xrightarrow{d_r}B_r\to B_{r+1}\to0$
is exact.

\emph{Step 2: exhaustion, before localization.} Let $y\in G^p$ represent a
class of $B_1^p$. By Proposition~\ref{paper:prop:x-surjective}, $y=zx_1$ for
\emph{one} $z\in Q$, and by exhaustiveness $z\in G^s$ for some $s$. For
$r\ge\max(1,p-s+1)$ we have $z\in G^{p-r+1}$. Hence
$y\in G^p\cap f(G^{p-r+1})$, and the class of $y$ vanishes in $B_r^p$. In
terms of order: if $y\in F_mQ$ and the chosen preimage lies in $F_{m'}Q$, the
class dies on page $r=\max(1,m'-m+1)$. The maps $\beta_r\colon B_1\to B_r$,
the identity on representatives, are surjective and $T$-linear, and their
kernels increase with $r$. An element of $B_1$ has finitely many homogeneous
components, so taking the largest $r$ shows that every $u\in B_1$ has some
$r$ with $\beta_r(u)=0$.

\emph{Step 3: localization.} Let $S=T\setminus\mathfrak q$. By
Lemma~\ref{paper:lem:tangential-finiteness}, $S^{-1}B_1\cong U_{\mathfrak q}$ has
finite length, so it is Noetherian over $T_{\mathfrak q}$. The maps
$S^{-1}\beta_r$ are surjective with increasing kernels. Each $u/s$ is killed
by $S^{-1}\beta_r$ for the $r$ belonging to $u$. The kernels therefore
stabilize at all of $S^{-1}B_1$, and $S^{-1}B_R=0$ for some $R$.

\emph{Step 4: telescoping.} Localizing the exact sequences
$0\to A_{r+1}\to A_r\xrightarrow{d_r}B_r\to B_{r+1}\to0$, all modules have
finite length because $A_1\cong V$ and $B_1\cong U$ do, and subsequent pages
are submodules and quotients. Additivity gives
$\ell(S^{-1}B_r)+\ell(S^{-1}A_{r+1})=
\ell(S^{-1}A_r)+\ell(S^{-1}B_{r+1})$.
Summing for $r=1,\ldots,R-1$ cancels every intermediate term and gives
$\ell(U_{\mathfrak q})+\ell(S^{-1}A_R)=\ell(V_{\mathfrak q})
+\ell(S^{-1}B_R)=\ell(V_{\mathfrak q})$.
\end{AIaddition}

\par\bigskip\noindent\pdfbookmark[1]{Comparison COMP-07}{bookmark-COMP-07}\hypertarget{comparison-COMP-07}{\textbf{Comparison COMP-07}}\label{comparison-COMP-07}\par\smallskip
\AIcomment{COMP-07}{\textbf{Information --- proof comparison; no edit requested.} The formal proof uses Gabber directly for minimal primes. If such a prime contained $x_1$, the $N$-fold bracket of the monic symbol $P$ with $x_1$ would put the unit $\pm N!$ in that prime. The subsequent stable-kernel length computation and Nakayama argument agree with the printed proof. \href{https://itpplasma.github.io/stafford38-formal/lean_proof_details.pdf\#nameddest=comparison-COMP-07}{Mathematical proof (COMP-07, p.~5).}}
\begin{AIaddition}{Proof as formalized in Lean}
The length computation (stable kernel $K=\ker x_1^m$, the split
$\ell(U_{\mathfrak q})=\ell(V_{\mathfrak q})+\ell((M'/K)/x_1(M'/K))$, and
Nakayama) is formalized as printed. Only the first step differs: Lean does
not pass through Theorem~\ref{paper:ithm:gabber} and
Lemma~\ref{paper:lem:componentwise-involutive}. Its formal Gabber theorem is stated
directly for minimal primes: every minimal prime $\mathfrak p$ of
$\Ann_{k[x,\xi]}\gr Q$ satisfies $\{\mathfrak p,\mathfrak p\}\subseteq\mathfrak p$.
Suppose $x_1\in\mathfrak p$. Since $P\in\Ann\gr Q\subseteq\mathfrak p$
(Corollary~\ref{paper:cor:charP}), induction gives
$\{x_1,\{x_1,\ldots,\{x_1,P\}\}\}\in\mathfrak p$ ($j$ brackets) for every $j$.
With the paper's bracket, $\{x_1,g\}=-\partial g/\partial\xi_1$. By
Lemma~\ref{paper:lem:order-symbol}, $P$ is homogeneous of degree $N$ with
coefficient $1$ at $\xi_1^N$. The $N$-fold bracket is therefore the constant
$(-1)^N\partial^N_{\xi_1}P=(-1)^NN!$, which is a unit in characteristic zero,
contradicting $\mathfrak p\neq(1)$. Lean's bracket is the negative of the
paper's, so there the constant is $N!$. Lean then transfers this avoidance
of $x_1$ to the minimal primes of the annihilator of $M'=M_{\mathfrak q}$,
using $\Ann_{C'}M'=S^{-1}\Ann_CM$ for the finite module $M$.
\end{AIaddition}

\par\bigskip\noindent\pdfbookmark[1]{Comparison COMP-08}{bookmark-COMP-08}\hypertarget{comparison-COMP-08}{\textbf{Comparison COMP-08}}\label{comparison-COMP-08}\par\smallskip
\AIcomment{COMP-08}{\textbf{Information --- proof comparison; no edit requested.} The formal conclusion is scheme-theoretic: $\Supp(\gr Q)\cap V(x_1)=\varnothing$, equivalently $\Ann(\gr Q)+(x_1)=(1)$. It therefore persists under field extension. The formal statement uses the support transported by $\xi\mapsto-\xi$; this fixes $x_1$, so the avoidance statements are equivalent. \href{https://itpplasma.github.io/stafford38-formal/lean_proof_details.pdf\#nameddest=comparison-COMP-08}{Mathematical proof (COMP-08, p.~5).}}
\begin{AIaddition}{Proof as formalized in Lean}
Lean states the conclusion scheme-theoretically:
$\Supp_{k[x,\xi]}(\gr Q)\cap V(x_1)=\varnothing$ in $\Spec k[x,\xi]$, for an
arbitrary field $k$ of characteristic zero. For the finite module $M=\gr Q$
this is equivalent to $\Ann(\gr Q)+(x_1)=(1)$. It is an identity of ideals, so
it persists under every field extension, and it gives the point-set statement
over any algebraically closed extension. The proof is the printed one:
Lemma~\ref{paper:lem:Uzero} gives $M/x_1M=0$, and
$\Supp(M/x_1M)=\Supp M\cap V(x_1)$ for finite $M$.

The formal statement is phrased for the transposed support, i.e.\ the support
of $\gr$ of the left module obtained from $Q$ through the transposition
$x_i\mapsto x_i$, $\partial_i\mapsto-\partial_i$. That support is the preimage
of $\Supp\gr Q$ under $\xi\mapsto-\xi$. This involution fixes $x_1$ and hence
$V(x_1)$, so the two disjointness statements are equivalent, and the formal
consumer converts back at once to $\Ann(\gr Q)+(x_1)=(1)$. No left-module
input is used anywhere.
\end{AIaddition}

\par\bigskip\noindent\pdfbookmark[1]{Comparison COMP-10}{bookmark-COMP-10}\hypertarget{comparison-COMP-10}{\textbf{Comparison COMP-10}}\label{comparison-COMP-10}\par\smallskip
\AIcomment{COMP-10}{\textbf{Information --- proof comparison; no edit requested.} The formal tangent-limit criterion is auxiliary. It assumes a generically smooth arc and a direct-summand lattice whose generic span is the lifted tangent space, with the specified axis component vanishing on its residue. Its proof does not use the smoothness hypothesis or the stated rank, but these remain assumptions of the formal statement. The formal proof of Theorem~\ref{paper:thm:asymptotic-conormal} follows a separate Laurent-series route, so it does not certify every step of the printed tangent-limit argument. \href{https://itpplasma.github.io/stafford38-formal/lean_proof_details.pdf\#nameddest=comparison-COMP-10}{Mathematical proof (COMP-10, p.~5).} ROUTE-01 states the correspondence status.}
\begin{AIaddition}{Formal status of Lemma~\ref{paper:lem:tangent-limit}}
Lean proves the following, over any algebraically closed field $K$ with no
characteristic assumption. Let $\mathfrak p\subset K[x_1,\dots,x_n]$ be prime,
$q\in K[[t]]^{n+1}$ with $q_0\neq0$, and fix an axis $1\le s\le n$. Assume:
\begin{itemize}[leftmargin=*]
\item some $q_i$ is a unit, and every homogenization $f^h$ with $f\in\mathfrak p$ satisfies $f^h(q)=0$ in $K((t))$;
\item $y=(q_i/q_0)_{i\ge1}$ is a smooth point of $\mathfrak p\otimes K((t))$;
\item $L\subset K[[t]]^{n+1}$ is a direct summand of rank $d+1$, with $d$ arbitrary;
\item the $K((t))$-span of $L$ equals the preimage of the Zariski tangent space of $\mathfrak p\otimes K((t))$ at $y$ under the derivative of dehomogenization at $q$;
\item every $w\in L$ has $w_s(0)=0$.
\end{itemize}
Then $[dx_s]$ lies in the Zariski closure of the smooth conormal directions of $V(\mathfrak p)$, that is, in $\Dir(V(\mathfrak p))$.
The smoothness hypothesis is carried but not needed for the conclusion, and the
rank $d$ plays no role. This criterion is auxiliary. The proof of
Theorem~\ref{paper:thm:asymptotic-conormal} does not use
Lemma~\ref{paper:lem:tangent-limit}; it uses the valuation-theoretic argument
given above after that theorem's proof.
\end{AIaddition}

\par\bigskip\noindent\pdfbookmark[1]{Comparison COMP-11}{bookmark-COMP-11}\hypertarget{comparison-COMP-11}{\textbf{Comparison COMP-11}}\label{comparison-COMP-11}\par\smallskip
\AIcomment{COMP-11}{\textbf{Checked endpoint.} Over an algebraically closed field of characteristic zero, the theorem proves $[dx_1]\in\operatorname{Dir}(Y)$. The constant-coordinate case is immediate. Otherwise one retained divisor valuation and its normalized projective column yield the residue-field tangent bound and the derivative columns used to construct a conormal covector with residue $e_1$. The equation-conormal closure gives the projective direction. The technical supplement gives the argument and states which steps of the printed proof remain unverified. \href{https://itpplasma.github.io/stafford38-formal/lean_proof_details.pdf\#nameddest=comparison-COMP-11}{COMP-11}.}
\begin{AIaddition}{Proof of Theorem~\ref{paper:thm:asymptotic-conormal} as formalized in Lean}
Let $K$ be algebraically closed of characteristic zero, let $\mathfrak p\subset
K[x_1,\dots,x_n]$ be prime, and suppose $Y=V(\mathfrak p)$ avoids $x_1=0$.
Write $R=K[x]/\mathfrak p$, $F=\operatorname{Frac}(R)$, and $\bar x_i$ for the
image of $x_i$. If $\bar x_1$ is algebraic over $K$, it is a nonzero constant;
then $x_1-c\in\mathfrak p$ for some $c\ne0$, and $e_1$ is the gradient of
this equation. So assume $\bar x_1$ is transcendental.

The construction chooses a divisorial valuation ring $V\subset F$ with
$\bar x_1\in\mathfrak m_V$, a uniformizer $t$, and residue field $\kappa$.
Completing $V$ gives a map into $\kappa[[t]]$ that sends the chosen uniformizer
to $t$. The same valuation witness supplies a normalized projective column
$q=(q_0,\dots,q_n)$, with some chart entry $q_j=1$ and
$q_1=q_0\bar x_1$. In the completion,
\[
 q_0=t^a u_0,\qquad q_1=t^b u_1,
 \qquad 0<a<b,\qquad u_0(0)\ne0.
\]
Dehomogenizing gives $y=(q_1/q_0,\dots,q_n/q_0)\in\kappa((t))^n$.
The induced field map $F\hookrightarrow\kappa((t))$ is injective, so for every
ground-field polynomial $f\in K[x_1,\dots,x_n]$, $f(y)=0$ exactly when
$f\in\mathfrak p$. Thus $y$ is the embedded generic point of the component,
with its full equation kernel.

The residues of the normalized column coordinates generate $\kappa$ up to an
algebraic extension. Since the characteristic is zero, that extension is
separable, and the differentials of those residues span $\Omega_{\kappa/K}$.
Set $d=\dim_\kappa\Omega_{\kappa/K}$. Choose $d$ residue coordinates whose
differentials form a basis, and choose the dual $K$-derivations
$D_1,\dots,D_d:\kappa\to\kappa$. Apply each derivation coefficientwise to
$q_i(t)=\sum_{r\ge0}c_{ir}t^r$ and set
$Z_{ij}(t)=\sum_{r\ge0}D_j(c_{ir})t^r$. On the chosen rows the constant
coefficient matrix of $Z$ is the identity. Its corresponding power-series
minor is therefore a unit.

The ratios $q_i/q_0$ generate $F$, and the same visible differential frame
bounds the component-field differentials.
The exact sequence
\[
 \mathfrak m_V/\mathfrak m_V^2\longrightarrow
 \kappa\otimes_V\Omega_{V/K}\longrightarrow\Omega_{\kappa/K}\longrightarrow0
\]
shows that the residue differential space needs $d$ generators and the
one-dimensional space $\mathfrak m_V/\mathfrak m_V^2$ contributes at most the
uniformizer direction. The finite visible differential module and Nakayama's
lemma then give
\[
 \dim_F\Omega_{F/K}\le d+1.
\]
The conormal sequence identifies $\Omega_{F/K}$ with the quotient of $F^n$
by the gradients of polynomials in $\mathfrak p$. After extending to
$\kappa((t))$, its dual is the tangent space at $y$. The full equation kernel
ensures these are exactly the equations of the component. Therefore
\[
\dim_{\kappa((t))}T_yV(\mathfrak p_{\kappa((t))})
 \le \dim_\kappa\Omega_{\kappa/K}+1.
\]

Now let $v=q'(t)$ be the uniformizer derivative. The columns $Z_{\bullet j}$
and $v$ live in projective coordinates; in the affine chart their quotient-rule
images have entries
\[
 \delta_j y_i=\frac{Z_{ij}q_0-q_iZ_{0j}}{q_0^2},\qquad
 \partial_t y_i=\frac{v_iq_0-q_iv_0}{q_0^2}\quad(1\le i\le n).
\]
Differentiating the base equations shows that these images are tangent vectors
at $y$. Since the selected minor of $Z$ is a unit, solve its $d$ selected
equations to choose coefficients
$\lambda_j(t)$ for which $v-Z\lambda$ vanishes on those rows. The row $q_0$
has order $a$, and every entry of $Z_{0\bullet}$ has order at least $a$.
But $q_0'=a t^{a-1}u_0+t^a u_0'$ has nonzero coefficient
$a u_0(0)$ in degree $a-1$. Characteristic zero therefore shows that the
correction cannot cancel that coefficient. If $c$ is the least order among
the entries of $v-Z\lambda$, then $c\le a-1$; divide by $t^c$ to obtain a
primitive vector $\tau$. Because $t$ is invertible in $\kappa((t))$, its
dehomogenized image is tangent as well.

In the $q_1$ row, both $q_1'$ and the coefficientwise columns have order at
least $b-1$. Since $c\le a-1<b-1$, the normalized entry $\tau_1$ still has
positive order, so $\tau_1(0)=0$. The selected entries of $\tau$ are zero by
the choice of $\lambda$. Form the matrix whose columns are the projective
position $q$, the $d$ coefficientwise columns $Z_{\bullet j}$, and $\tau$.
The chart row isolates $q$; the selected minor isolates the $Z$ columns; and
primitivity gives a nonzero pivot for $\tau$. Its reduction modulo $t$ has
full column rank, hence the power-series matrix has a left inverse.

At $t=0$ the distinguished $q_1$ row vanishes on every column: $q_1(0)=0$,
the $Z$ entries have order at least $b$, and $\tau_1(0)=0$. Thus the pure
axis covector annihilates the reduced matrix. The left inverse lifts it to an
exact power-series covector $\ell(t)$ annihilating the whole matrix, with
residue $\ell(0)$ equal to that axis covector. The quotient-rule formulas show
that each raw column $Z_{\bullet j}$ and $\tau$ lifts a tangent vector at $y$,
while $q$ is the radial position vector. Thus the $d+2$ independent columns
lie in the affine cone over the projective tangent plane at $[q]$. The affine
tangent at $y$ has dimension at most $d+1$ by the bound above, so this cone
has dimension at most $d+2$; the independent columns therefore span it. Since
$q_0\ne0$ in $\kappa((t))$, dehomogenization identifies the cone modulo its
radial line $\kappa((t))q$ with the affine tangent at $y$. The remaining
$d+1$ quotient-rule columns give a basis of that tangent space. By the full
equation kernel, the affine tail of $\ell(t)$ lies
in the equation-conormal locus for $\mathfrak p$ over $\kappa((t))$, and its
residue is $e_1$.

Finally, a homogeneous polynomial vanishing on the smooth conormal directions
vanishes on the equation-conormal locus: write a covector as a finite sum of
gradients, use density of the smooth locus and the Nullstellensatz, then extend
the resulting polynomial identity to $\kappa((t))$. Specializing along the
power-series covector shows that the polynomial vanishes at $e_1$. Hence
$[dx_1]\in\operatorname{Dir}(Y)$.

This proves the endpoint by a generic-divisor route. It does not identify the
retained valuation ring with the local ring of the divisor on the normalized
projective closure, or construct the paper's selected smooth closed point,
completed local coordinates, tilted arc, or tangent-frame specialization.
Those steps of the printed proof remain a separate correspondence question.
\end{AIaddition}

\par\bigskip\noindent\pdfbookmark[1]{Comparison COMP-13}{bookmark-COMP-13}\hypertarget{comparison-COMP-13}{\textbf{Comparison COMP-13}}\label{comparison-COMP-13}\par\smallskip
\AIcomment{COMP-13}{\textbf{Information --- no further paper edit requested.} The formal proof uses the same finite Hamiltonian translations and component conormal containment. The equality of component ideals on a dense open set is checked through a separator and the product rule, then extended by closure. The limiting direction is used in the closure of fibre directions, as in the marked GEO-03 correction. \href{https://itpplasma.github.io/stafford38-formal/lean_proof_details.pdf\#nameddest=comparison-COMP-13}{Mathematical proof (COMP-13, p.~8).}}
\begin{AIaddition}{Proof of Theorem~\ref{paper:thm:coisotropic-exclusion} as formalized in Lean}
Work in the setting of the formal statement and suppose
$W\cap\pi^{-1}(H)=\varnothing$. Put $J=I(W)$, which is radical, and $B=J\cap K[x]$.
\begin{enumerate}[label=\textbf{Step \arabic*.},ref=\arabic*,leftmargin=*,widest=5]
\item\label{st:hom} (\emph{$\xi$-homogeneity.}) Write $g\in J$ as $g=\sum_dg_d$ with $g_d$ homogeneous of
degree $d$ in $\xi$. For $(y,\xi)\in W$, the polynomial $a\mapsto\sum_da^dg_d(y,\xi)$
vanishes on the infinite set $K^\times$. So every $g_d\in J$, i.e.\ $J$ is $\xi$-homogeneous.
\item\label{st:zero} (\emph{Zero section.}) Let $Y=V(B)$. For $y\in Y$ and $g\in J$,
$g(y,0)=g_0(y)=0$ because $g_0\in J$ has $\xi$-degree $0$, i.e.\ $g_0\in B$. So
$Y\times\{0\}\subset V(J)=W$. Hence $\pi(W)=Y$ is closed, and
$Y\cap H=\varnothing$.
\item\label{st:transl} (\emph{Vertical translation.}) Let $f\in B$. Then $f\circ\pi\in J$, so
$H_fg:=\sum_i\partial_{x_i}f\,\partial_{\xi_i}g=\pm\{f,g\}$ lies in $J$ for every $g\in J$.
This is the only use of involutivity. Fix $y\in Y$ and $g\in J$ and let
$\varphi(t)=g(y,t\,df_y)\in K[t]$. The chain rule gives
$\varphi^{(m)}(t)=(H_f^mg)(y,t\,df_y)$, and $H_f^mg\in J$, so
$\varphi^{(m)}(0)=0$ by Step~\ref{st:zero}. As $m!\neq0$ in $K$, we get $\varphi=0$,
so $(y,t\,df_y)\in W$ for all $t$. Since
$\sum_ac_a\,d(f_a)_y=d(\sum_ac_af_a)_y$ and $\sum_ac_af_a\in B$, it follows that
$\{y\}\times\operatorname{span}\{df_y:f\in B\}\subset W$. Each translation is the finite polynomial Hamiltonian flow; the proof composes
these flows for the finite conormal decomposition, so no infinite exponential series is used.
\item\label{st:comp} (\emph{Component containment.}) Let $\mathfrak q$ be a minimal prime of $B$ and
$Y_0=V(\mathfrak q)$. Taking a product of elements of the other minimal primes
that lie outside $\mathfrak q$ gives $s\notin\mathfrak q$ with $s\mathfrak q\subset B$.
For $y\in Y_0$ with $s(y)\neq0$ and $h\in\mathfrak q$, we have $d(sh)_y=s(y)\,dh_y$, so
$(y,dh_y)\in W$ by Step~\ref{st:transl}. For $g\in J$, the polynomial
$G(x)=g(x,\nabla h(x))$ therefore vanishes on $Y_0\cap\{s\neq0\}$. Hence
$sG\in\mathfrak q$, so $G\in\mathfrak q$, and $(y,\nabla h(y))\in W$ for \emph{every}
$y\in Y_0$. Absorbing linear combinations into $h$ gives
$\{y\}\times\operatorname{span}\{\nabla h(y):h\in\mathfrak q\}\subset W$ for all $y\in Y_0$. In particular
$N^*_yY_0\subset W_y$ at every smooth point $y$ of $Y_0$, which proves conormal containment.
\item\label{st:asym} (\emph{Conclusion in the affine cone.}) $Y_0$ is irreducible, nonempty and
disjoint from $H$. By Step~\ref{st:comp} and $W\subset V(P)$, $P$ vanishes on
$N(Y_0)=\bigcup_{y\in(Y_0)_{\sm}}N^*_yY_0\subset K^n$. By part (i) of
Theorem~\ref{paper:thm:asymptotic-conormal} in its formal version, $P(1,0,\dots,0)=0$,
a contradiction.
\end{enumerate}
Part (i) is a statement about the vanishing ideal of the cone
$N(Y_0)\subset K^n$. So $[dx_1]$ never has to be a fibre direction of $W$ over a
finite base point, which removes the closure issue of GEO-03. Homogeneity of $P$ is
not used, and one component $Y_0$ suffices.
\end{AIaddition}

\par\bigskip\noindent\pdfbookmark[1]{Comparison COMP-15}{bookmark-COMP-15}\hypertarget{comparison-COMP-15}{\textbf{Comparison COMP-15}}\label{comparison-COMP-15}\par\smallskip
\AIcomment{COMP-15}{\textbf{Information --- no further paper edit requested.} Lean constructs the quotient scalar-extension isomorphism, checks the associated-graded annihilator comparison, and descends quotient vanishing by faithful scalar extension. Geometry runs over $\Kbar$; rank-zero and scalar cases are handled separately. \href{https://itpplasma.github.io/stafford38-formal/lean_proof_details.pdf\#nameddest=comparison-COMP-15}{Mathematical proof (COMP-15, p.~9).}}
\begin{AIaddition}{Proof as formalized in Lean}
Lean works over $k$ and runs the geometry over $\Kbar$. It checks the
paper's quotient scalar-extension and faithful-descent route. The proof has
six steps.

\emph{Step 1: degenerate cases.} For $n=0$, $A_0(k)=k$ and $d\neq0$ is a unit,
so $R=d^{-1}$ and $F=S=0$ work. Let $n\ge1$ and $N=\deg_Bd$. If $N=0$, then
$d=c\in k^\times$, and $\ell=x_1$, $F=\ell^0=1$, $R=c^{-1}$, $S=0$ work.

\emph{Step 2: normal form.} For $N\ge1$, the formal version of
Proposition~\ref{paper:prop:monic} gives symplectic $M,M'$ with
$MM'=M'M=1$ and $c\in k^\times$ such that
$\tilde d=c^{-1}\Phi_M(d)\in A_n(k)$ has Bernstein degree at most $N$ and
coefficient $1$ on $p^N$. Let
$I=\tilde dA_n+x_1^N\tilde dA_n$ and $Q=A_n/I$. It suffices to show $Q=0$.

\emph{Step 3: vanishing over $\Kbar$.} The image of $\tilde d$ in
$A_n(\Kbar)$ is again PBW-monic, so Sections 5--7 apply to
$Q_{\Kbar}=A_n(\Kbar)/I_{\Kbar}$. Suppose
$J:=\sqrt{\operatorname{in}(I_{\Kbar})}\neq(1)$, and let
$W=V(J)\subseteq\Kbar^{2n}$. Then:
\begin{itemize}[nosep,leftmargin=1.4em]
\item $W\neq\varnothing$ by the Nullstellensatz, and $I(W)=J$.
\item $W$ is closed and fibre-conical. The generators $\sigma_L(z)$ of
$\operatorname{in}(I_{\Kbar})$ are homogeneous of degree $L$ in $\xi$, so
$\operatorname{in}(I_{\Kbar})$ and $J$ are $\xi$-homogeneous. For $f\in J$
each $\xi$-homogeneous component $f_j$ lies in $J$, and
$f(y,t\xi)=\sum_jt^jf_j(y,\xi)$.
\item $W$ is involutive: $J=\sqrt{\Ann\gr Q_{\Kbar}}$ is involutive by
Theorem~\ref{paper:ithm:gabber} (radical form, a corollary of the minimal-prime
form).
\item $P=\sigma(\tilde d)\in\operatorname{in}(I_{\Kbar})\subseteq J$, so
$W\subseteq V(P)$, and $P(1,0,\ldots,0)=1$ by Lemma~\ref{paper:lem:order-symbol}.
\item $W\cap\pi^{-1}(H)=\varnothing$ by Proposition~\ref{paper:prop:char-avoids}
over $\Kbar$, i.e.\ $\operatorname{in}(I_{\Kbar})+(x_1)=(1)$.
\end{itemize}
The formal Theorem~\ref{paper:thm:coisotropic-exclusion} gives a point of $W$ with
$x_1=0$, a contradiction. Hence $\operatorname{in}(I_{\Kbar})=(1)$.

\emph{Step 4: quotient and descent.} Lean constructs
$Q\otimes_k\Kbar\cong A_n(\Kbar)/I_{\Kbar}$ using the actual quotient
and coefficient-extension maps. It also proves
$\Ann\gr Q_{\Kbar}=(\Ann\gr Q)\Kbar[x,\xi]$.
Empty characteristic support gives $Q_{\Kbar}=0$ by the exhaustive
filtration. Faithful scalar extension of vector spaces then gives $Q=0$.

\emph{Step 5: the certificate over $k$.} Since $Q=0$, we have $1\in I$.
So $1=\tilde dR+x_1^N\tilde dS$ with $R,S\in A_n(k)$.

\emph{Step 6: transport back.} Since $c^{-1}$ is central, we get
\[1=\Phi_M(d)(c^{-1}R)+x_1^N\Phi_M(d)(c^{-1}S).\] Applying
$\Phi_M^{-1}=\Phi_{M'}$ gives $1=dR'+\ell^NdS'$ with
$R'=\Phi_{M'}(c^{-1}R)$, $S'=\Phi_{M'}(c^{-1}S)$ and
$\ell=\Phi_{M'}(x_1)=\sum_jM'_{1j}z_j$, where row $1$ is the row of $x_1$
and $N=\deg_Bd$.
\end{AIaddition}

\par\bigskip\noindent\pdfbookmark[1]{Comparison COMP-16}{bookmark-COMP-16}\hypertarget{comparison-COMP-16}{\textbf{Comparison COMP-16}}\label{comparison-COMP-16}\par\smallskip
\begin{AIaddition}{Proof as formalized in Lean}
No two-generator theorem and no Ore condition are needed. Let $T$ be a finitely
generated torsion right module over a ring $A$ without zero divisors in which
every $d\ne0$ admits $1=dR+FdS$; $A=A_n(k)$ qualifies by Theorem~\ref{paper:thm:main}
and because $A_n(k)$ is a domain (the Bernstein symbol of a product is the
product of the symbols, in the polynomial ring $k[x,\xi]$).

\emph{The paper's matrix.} Let $t_1,t_2\in T$. Choose $d_2\ne0$ with $t_2d_2=0$ and
$e\ne0$ with $(t_1d_2)e=0$; then $d:=d_2e\ne0$ annihilates $t_1$ and $t_2$. Write
$1=dR+FdS$ and put $a=dS$. The two columns $(a,1-Fa)$ and $(1,-F)$
span $A\oplus A$: for $(x,y)$ take $u=y+Fx$ and $v=x-a(y+Fx)$.
The first column lies in the kernel of $(r,s)\mapsto t_1r+t_2s$.
Thus its image is generated by $g=t_1-t_2F$, and
$t_1A+t_2A=gA$. This is the displayed matrix calculation in the paper,
checked directly.

\emph{Induction.} If $t_1,\dots,t_m$ generate $T$, replace $t_{m-1},t_m$ by a
single generator of $t_{m-1}A+t_mA$ and repeat; after $m-1$ steps one generator
remains.
\end{AIaddition}

\ifshowleandetails
\section*{Declaration links}
\AIadd{\formalref[{Lean also covers $n=0$, where $A_0(k)=k$}]{%
\leandecl{Stafford38/FoundationClosure.lean}{67}{Stafford38.universalStatement} proves
\leandecl{Stafford38/Statement.lean}{23}{Stafford38.UniversalStatement};
\leandecl{Stafford38/FoundationClosure.lean}{71}{Stafford38.universalFixedSourceStatement} proves
\leandecl{Stafford38/FixedSourceStatement.lean}{35}{Stafford38.FixedSource.UniversalFixedSourceStatement}
with \leandecl{Stafford38/FixedSourceStatement.lean}{28}{Stafford38.FixedSource.IsLinearWeylCoordinate}
and \leandecl{Stafford38/FixedSourceStatement.lean}{23}{Stafford38.FixedSource.bernsteinDegree};
Mathlib-only statements compared by Comparator:
\leandecl{Stafford38/ChallengeDefinitions.lean}{78}{Stafford38Challenge.UniversalStatement},
\leandecl{Stafford38/ChallengeDefinitions.lean}{193}{Stafford38FixedSourceChallenge.UniversalFixedSourceStatement}.}}

\AIadd{\formalref{%
\leandecl{Stafford38/Weyl/IteratedEquivalence.lean}{636}{Stafford38.WeylIteratedEquivalence.presentedIteratedEquiv},
\leandecl{Stafford38/Weyl/OuterOreMonic.lean}{47}{Stafford38.WeylOuterOreMonic.presentedOuterPolynomial}}}

\AIadd{\formalref{\leandecl{Stafford38/Characteristic/Polynomial.lean}{22}{Stafford38.Characteristic.SymbolRing}
(the ring $k[x,\xi]$, with its base and cotangent coordinate blocks $x_i$ and $\xi_i$)}}

\AIadd{\formalref[{as spans of ordered PBW words of bounded weight}]{%
\leandecl{Stafford38/Weyl/Filtration.lean}{284}{Stafford38.WeylFiltration.bernsteinPiece},
\leandecl{Stafford38/Weyl/Filtration.lean}{288}{Stafford38.WeylFiltration.orderPiece},
\leandecl{Stafford38/FixedSourceStatement.lean}{23}{Stafford38.FixedSource.bernsteinDegree}}}

\AIadd{\formalref[{for $M=A_n/I$}]{%
\leandecl{Stafford38/Characteristic/InitialIdeal.lean}{189}{Stafford38.CharacteristicInitialIdeal.orderInitialIdeal},
\leandecl{Stafford38/Characteristic/InitialIdeal.lean}{197}{Stafford38.CharacteristicInitialIdeal.orderCharacteristicSupport},
\leandecl{Stafford38/Characteristic/AssociatedGradedModule.lean}{373}{Stafford38.CharacteristicAssociatedGradedModule.annihilator_orderAssociatedGradedModule}}}

\AIadd{\formalref[{Lean's bracket is the negative of this one; involutivity is unaffected}]{%
\leandecl{Stafford38/Characteristic/Polynomial.lean}{31}{Stafford38.Characteristic.poissonBracket},
\leandecl{Stafford38/Characteristic/SymplecticCompletion.lean}{31}{Stafford38.CharacteristicSymplecticCompletion.phasePairing}}}

\AIadd{\formalref{\leandecl{Stafford38/Characteristic/PostScalarExtensionPoisson.lean}{33}{Stafford38.Characteristic.PostScalarExtensionPoisson.IsInvolutive};
negative control \leandecl{Stafford38/Geometry/GeneralCoisotropicSetsTest.lean}{164}{Stafford38.Geometry.GeneralCoisotropicSetsTest.zeroSectionIdealOne_not_isPoisson}}}

\AIadd{\formalref[{proved in Lean, not cited; exactly this scope, over every field of characteristic zero}]{%
\leandecl{Stafford38/Characteristic/GabberGlobalAssembly.lean}{224}{Stafford38.Characteristic.GabberGlobalAssembly.weylAssociatedGradedRadicalInvolutivity},
statement \leandecl{Stafford38/Characteristic/CanonicalGabberInvolutivityInterface.lean}{86}{Stafford38.Characteristic.CanonicalGabberInvolutivityInterface.WeylAssociatedGradedRadicalInvolutivity},
minimal primes \leandecl{Stafford38/Characteristic/GabberGlobalAssembly.lean}{144}{Stafford38.Characteristic.GabberGlobalAssembly.minimalPrime_isInvolutive}}}

\AIadd{\formalref{%
\leandecl{Stafford38/Characteristic/PaperLocalizedInvolutivity.lean}{187}{Stafford38.Characteristic.PaperLocalizedInvolutivity.minimalPrimes_isInvolutive_of_radical_isInvolutive_viaLocalization}}}

\AIadd{\formalref{\leandecl{Stafford38/Weyl/Symplectic.lean}{46}{Stafford38.WeylSymplectic.standardSymplecticAlgHom},
\leandecl{Stafford38/Weyl/Symplectic.lean}{91}{Stafford38.WeylSymplectic.standardSymplecticAlgEquivOfInverse}
(with explicit inverse), \leandecl{Stafford38/Weyl/SymbolCompatibility.lean}{416}{Stafford38.WeylSymbolCompatibility.standardSymplecticAlgHom_principal_compatibility}}}

\AIadd{\formalref[{completion to a symplectic basis, then the basis-to-basis isometry}]{%
\leandecl{Stafford38/Characteristic/SymplecticCompletion.lean}{238}{Stafford38.CharacteristicSymplecticCompletion.exists_symplectic_mulVec_eq}}}

\AIadd{\formalref[{monicity is stated as: $\tilde d\in B_N$ and the PBW coefficient of $\partial_1^N$ is $1$}]{%
\leandecl{Stafford38/Weyl/MonicNormalization.lean}{158}{Stafford38.WeylMonicNormalization.exists_normalized_symplectic_image},
\leandecl{Stafford38/Weyl/MonicNormalization.lean}{181}{Stafford38.WeylMonicNormalization.scalar_or_normalized_symplectic_image},
\leandecl{Stafford38/Weyl/PBWMonicBridge.lean}{98}{Stafford38.WeylPBWMonicBridge.IsPBWMonicAt},
\leandecl{Stafford38/Characteristic/SymplecticCompletion.lean}{346}{Stafford38.CharacteristicSymplecticCompletion.exists_symplectic_chart_matrices}}}

\AIadd{\formalref{\leandecl{Stafford38/UniversalAssembly.lean}{69}{Stafford38.UniversalAssembly.universalStatement_of_canonicalSupportVanishing},
\leandecl{Stafford38/FixedSourceAssembly.lean}{24}{Stafford38.FixedSource.universalFixedSourceStatement_of_canonicalSupportVanishing}}}

\AIadd{\formalref{\texttt{Polynomial.Monic} after \leandecl{Stafford38/Characteristic/NormalSymbolPolynomial.lean}{28}{Stafford38.Characteristic.NormalSymbolPolynomial.normalSymbolAlgEquiv}}}

\AIadd{\formalref{\leandecl{Stafford38/Characteristic/InitialIdeal.lean}{320}{Stafford38.CharacteristicInitialIdeal.canonical_orderPrincipalComponent_isFibreOnly},
\leandecl{Stafford38/Geometry/ConormalAxisContradiction.lean}{61}{Stafford38.Geometry.ConormalAxisContradiction.canonical_orderPrincipalComponent_isHomogeneous},
\leandecl{Stafford38/Characteristic/NormalSymbolPolynomial.lean}{104}{Stafford38.Characteristic.NormalSymbolPolynomial.canonicalNormalPolynomial_monic},
\leandecl{Stafford38/Geometry/ConormalAxisContradiction.lean}{131}{Stafford38.Geometry.ConormalAxisContradiction.canonical_orderPrincipalComponent_eval_pureMomentumAxis}}}

\AIadd{\formalref{\leandecl{Stafford38/Weyl/EulerResidue.lean}{26}{Stafford38.WeylEulerResidue.canonicalRightIdeal}
$=\operatorname{span}\{d,x^Nd\}$ as a right ideal}}

\AIadd{\formalref[{Lean uses the subring generated by the transverse algebra, $x_1$ and $x_1p$, and one-sided normality $x_1r=r'x_1$, which is all that is used}]{%
\leandecl{Stafford38/Weyl/EulerSubring.lean}{25}{Stafford38.WeylEulerSubring.eulerSubring},
\leandecl{Stafford38/Weyl/EulerSubring.lean}{71}{Stafford38.WeylEulerSubring.coordinate_normal_eulerSubring}}}

\AIadd{\formalref{\leandecl{Stafford38/Weyl/EulerProductIdentities.lean}{70}{Stafford38.Weyl.EulerProductIdentities.eval_fallingEulerProduct} and
\leandecl{Stafford38/Weyl/EulerProductIdentities.lean}{79}{Stafford38.Weyl.EulerProductIdentities.eval_risingEulerProduct} (the corrected products),
\leandecl{proofs/weyl_pure_power.lean}{148}{Stafford.exists_eulerPolynomial_x_pow_mul_d_pow}}}

\AIadd{\formalref[{the ring identity is Lean's; the factorization uses left coefficients $c\,x^ap^j$}]{%
\leandecl{Stafford38/Weyl/EulerResidue.lean}{41}{Stafford38.WeylEulerResidue.positiveEulerResidue_identity},
\leandecl{Stafford38/Weyl/EulerRemainder.lean}{257}{Stafford38.WeylEulerRemainder.positiveEulerResidue_eq_eulerSubring_mul_coordinate},
\leandecl{Stafford38/Weyl/EulerRemainder.lean}{405}{Stafford38.WeylEulerRemainder.presentedPositiveEulerResidue}}}

\AIadd{\formalref{\leandecl{Stafford38/Weyl/QuotientTransport.lean}{240}{Stafford38.WeylQuotientTransport.presentedCanonicalRightQuotient_rightMul_coordinate_surjective},
\leandecl{Stafford38/Quotient/EulerSurjectivity.lean}{86}{Stafford38.EulerSurjectivity.rightMul_surjective_of_euler_normal};
generation by $C$ and $p$: \leandecl{Stafford38/Weyl/EulerRemainder.lean}{310}{Stafford38.WeylEulerRemainder.closure_pairOldSubring_union_coordinate_momentum_eq_top}}}

\AIadd{\formalref{\texttt{Module.length} (Mathlib); finite length of the localized modules:
\leanlib{AlgebraicAnalysis/Module/MinimalSupportKernelCokernelLengths.lean}{48}{AlgebraicAnalysis.localized_kernel_and_cokernel_isFiniteLength}}}

\AIadd{\formalref{\leandecl{Stafford38/Characteristic/CanonicalOldTangentialFiniteness.lean}{73}{Stafford38.Characteristic.CanonicalOldTangentialFiniteness.canonical_finite_old_coordinate_kernel_cokernel},
\leandecl{Stafford38/Characteristic/CanonicalTangentialSymbolFiniteness.lean}{116}{Stafford38.Characteristic.CanonicalTangentialSymbolFiniteness.canonical_orderAssociatedGradedModule_finite_tangential_coordinate}}}

\AIadd{\formalref{\leandecl{Stafford38/Characteristic/CanonicalPageEulerInequality.lean}{25}{Stafford38.Characteristic.CanonicalPageEulerInequality.canonicalPage_length_target_le_source};
exhaustion \leanlib{AlgebraicAnalysis/Module/UniformBoundaryVanishing.lean}{106}{AlgebraicAnalysis.exists_uniform_subsingleton_localized}}}

\AIadd{\formalref{\leanlib{AlgebraicAnalysis/Module/BaseLocalizedKoszulPositivity.lean}{25}{AlgebraicAnalysis.BaseLocalizedKoszulPositivity.localized_length_cokernel_gt_kernel};
$x_1$ in no minimal prime: \leandecl{Stafford38/Characteristic/NoncharacteristicMinimalPrime.lean}{118}{Stafford38.Characteristic.NoncharacteristicMinimalPrime.canonical_minimalPrime_mem_of_normalCoordinate_false}}}

\AIadd{\formalref{\leandecl{Stafford38/Characteristic/CanonicalKoszulContradiction.lean}{86}{Stafford38.Characteristic.CanonicalKoszulContradiction.coordinate_cokernel_subsingleton},
\leandecl{Stafford38/Characteristic/CanonicalKoszulContradiction.lean}{50}{Stafford38.Characteristic.CanonicalKoszulContradiction.no_minimal_coordinate_cokernel_support}}}

\AIadd{\formalref{\leandecl{Stafford38/Characteristic/InitialIdeal.lean}{281}{Stafford38.CharacteristicInitialIdeal.canonical_orderPrincipalComponent_mem_initialIdeal},
\leandecl{Stafford38/Characteristic/InitialIdeal.lean}{291}{Stafford38.CharacteristicInitialIdeal.canonical_orderCharacteristicSupport_subset_principal_zeroLocus}}}

\AIadd{\formalref[{stated for the support transported along $\xi\mapsto-\xi$, which fixes $x_1$; transported back verbatim}]{%
\leandecl{Stafford38/Characteristic/CanonicalKoszulContradiction.lean}{123}{Stafford38.Characteristic.CanonicalKoszulContradiction.canonical_support_avoidance},
\leandecl{Stafford38/Characteristic/SpecializedNoncharacteristicEquality.lean}{151}{Stafford38.SpecializedNoncharacteristicEquality.transposedSupport_disjoint_axis_iff}}}

\AIadd{\formalref[{$Y$ given by a prime ideal; conormal space = annihilator of the Zariski tangent space; closure by homogeneous polynomials}]{%
\leandecl{Stafford38/Geometry/ProjectiveConormalDirections.lean}{40}{Stafford38.Geometry.ProjectiveConormalDirections.smoothConormalDirectionSet},
\leandecl{Stafford38/Geometry/ProjectiveConormalDirections.lean}{32}{Stafford38.Geometry.ProjectiveConormalDirections.projectiveHomogeneousClosure}}}

\AIadd{\formalref[{any algebraically closed field and any axis; auxiliary, not used by the formal proof of Theorem~\ref{paper:thm:asymptotic-conormal}}]{%
\leandecl{Stafford38/Geometry/GeneralTangentLimitCriterion.lean}{629}{Stafford38.Geometry.GeneralTangentLimitCriterion.tangent_limit_criterion_of_directSummand};
scalar extension \leandecl{Stafford38/Geometry/ConormalScalarExtensionVanishing.lean}{137}{Stafford38.Geometry.ConormalScalarExtensionVanishing.scalarExtension_vanishing}}}

\AIadd{\formalref[{generic point only Zariski-tangent; no smoothness needed for this step}]{%
\leandecl{Stafford38/Geometry/ConormalScalarExtensionVanishing.lean}{137}{Stafford38.Geometry.ConormalScalarExtensionVanishing.scalarExtension_vanishing};
\leandecl{Stafford38/Geometry/SmoothConormalFibreVanishing.lean}{24}{Stafford38.Geometry.SmoothConormalFibreVanishing.fibreLift_mem_vanishingIdeal_equationConormal};
\leandecl{Stafford38/Geometry/LaurentConormalDirection.lean}{71}{Stafford38.Geometry.LaurentConormalDirection.residue_mem_extensionFibreClosure_of_laurent_generic};
\leandecl{Stafford38/Geometry/GeneralTangentLimitCriterion.lean}{501}{Stafford38.Geometry.GeneralTangentLimitCriterion.exists_axis_laurent_smooth_conormal_direction}}}

\AIadd{\formalref[{over every algebraically closed field of characteristic zero; the existing Laurent proof has the same conclusion, while its intermediate route differs from the printed proof}]{%
\leandecl{Stafford38/Geometry/GeneralAsymptoticConormal.lean}{42}{Stafford38.Geometry.GeneralAsymptoticConormal.coordinate_axis_mem_projective_conormal_directions},
affine-cone form \leandecl{Stafford38/Geometry/GeneralAsymptoticConormal.lean}{32}{Stafford38.Geometry.GeneralAsymptoticConormal.coordinate_axis_mem_smooth_fibre_closure}}}

\AIadd{\formalref[{any algebraically closed field of characteristic $0$; (i) is the affine-cone form used in Section~\ref{paper:sec:involutive-review}}]{%
\leandecl{Stafford38/Geometry/GeneralAsymptoticConormal.lean}{42}{Stafford38.Geometry.GeneralAsymptoticConormal.coordinate_axis_mem_projective_conormal_directions};
\leandecl{Stafford38/Geometry/GeneralDivisorialVisibleFrame.lean}{339}{Stafford38.Geometry.GeneralDivisorialVisibleFrame.generalDivisorialVisibleFrameExistence};
\leandecl{Stafford38/Geometry/DivisorTangentLattice.lean}{335}{Stafford38.Geometry.DivisorTangentLattice.VisibleDivisorFrame.exists_maximalIdeal_coefficients};
\leandecl{Stafford38/Geometry/ConormalScalarExtensionVanishing.lean}{137}{Stafford38.Geometry.ConormalScalarExtensionVanishing.scalarExtension_vanishing}}}

\AIadd{\formalref[{any algebraically closed field of characteristic $0$; $P$ need not be homogeneous}]{%
\leandecl{Stafford38/Geometry/GeneralCoisotropicSets.lean}{160}{Stafford38.Geometry.GeneralCoisotropicSets.exists_zero_base_coordinate_of_isFibreConical};
\leandecl{Stafford38/Geometry/GeneralCoisotropicExclusion.lean}{84}{Stafford38.Geometry.GeneralCoisotropicExclusion.exists_zero_base_coordinate};
\leandecl{Stafford38/Characteristic/BaseRelativePoisson.lean}{29}{Stafford38.Characteristic.BaseRelativePoisson.IsBaseRelativePoisson};
\leandecl{Stafford38/Geometry/GeneralCoisotropicSetsTest.lean}{104}{Stafford38.Geometry.GeneralCoisotropicSetsTest.exact_complex_manuscript_coisotropic_consumer}}}

\AIadd{\formalref{%
\leandecl{Stafford38/Geometry/GeneralCoisotropicExclusion.lean}{84}{Stafford38.Geometry.GeneralCoisotropicExclusion.exists_zero_base_coordinate};
\leandecl{Stafford38/Characteristic/BaseRelativePoisson.lean}{72}{Stafford38.Characteristic.BaseRelativePoisson.zeroSection_stable_under_differential_translation_of_isBaseRelativePoisson};
\leandecl{Stafford38/Geometry/GeneralComponentConormalContainment.lean}{72}{Stafford38.Geometry.GeneralComponentConormalContainment.equationConormalClosure_minimalPrime_subset_zeroLocus};
\leandecl{Stafford38/Geometry/GeneralAsymptoticConormal.lean}{32}{Stafford38.Geometry.GeneralAsymptoticConormal.coordinate_axis_mem_smooth_fibre_closure}}}

\AIadd{\formalref{\leandecl{Stafford38/FoundationClosure.lean}{41}{Stafford38.FoundationClosure.canonicalSupportVanishingViaGeneralCoisotropic}:
exclusion over $\overline k$ \leandecl{Stafford38/Geometry/GeneralCoisotropicCanonicalAdapter.lean}{222}{Stafford38.Geometry.GeneralCoisotropicCanonicalAdapter.algebraicallyClosedCanonicalSupportVanishing_of_generalCoisotropic},
descent \leandecl{Stafford38/Weyl/FilteredScalarLifting.lean}{472}{Stafford38.Weyl.FilteredScalarLifting.canonicalSupportDescent},
certificate \leandecl{Stafford38/Characteristic/CanonicalCertificate.lean}{33}{Stafford38.CharacteristicCanonicalCertificate.exists_fixedSource_certificate_of_orderCharacteristicSupport_eq_empty}}}

\AIadd{\formalref[{stated for the canonical ideal; the argument uses nothing else}]{%
\leandecl{Stafford38/Weyl/FilteredScalarLifting.lean}{387}{Stafford38.Weyl.FilteredScalarLifting.target_orderInitialIdeal_le_map_source_orderInitialIdeal};
\leandecl{Stafford38/Weyl/PresentedScalarExtension.lean}{390}{Stafford38.Weyl.PresentedScalarExtension.presentedWeylScalarExtension_map_orderInitialIdeal_le};
\leandecl{Stafford38/Characteristic/EmptySupportVanishing.lean}{83}{Stafford38.CharacteristicEmptySupportVanishing.orderInitialIdeal_eq_top_iff_rightQuotient_subsingleton};
\leandecl{Stafford38/Characteristic/GeometricSupportDescent.lean}{101}{Stafford38.Characteristic.GeometricSupportDescent.map_scalarPolynomialMap_eq_top_iff}}}

\AIadd{\formalref{\leandecl{Stafford38/TorsionCyclicity.lean}{66}{Stafford38.TorsionCyclicity.weyl_isCyclic_of_isRightTorsion},
two generators \leandecl{Stafford38/TorsionCyclicity.lean}{52}{Stafford38.TorsionCyclicity.exists_span_singleton_eq_span_pair},
domain \leandecl{Stafford38/Weyl/Domain.lean}{25}{Stafford38.WeylDomain.mul_ne_zero};
Mathlib-only statement \leandecl{CorollaryChallenge.lean}{34}{Stafford38CorollaryChallenge.torsionCyclicStatement}}}

\AIadd{\formalref[{for $X=\mathbb A^n$, $H=\{x_1=0\}$ and a characteristic ideal $J$: the coordinate-ring map $k[T^*H]\to k[x,\xi]/(J+(x_1))$ is finite}]{%
\leandecl{Stafford38/NoncharacteristicHyperplane.lean}{73}{Stafford38.NoncharacteristicHyperplane.IsNoncharacteristic}}}

\AIadd{\formalref[{the zero-section statement holds without $x_1=0$: tangential covariables vanishing forces $\xi_1=0$}]{%
\leandecl{Stafford38/NoncharacteristicHypersurface.lean}{132}{Stafford38.NoncharacteristicHypersurface.canonical_principal_hypersurface_finite}, \leandecl{Stafford38/NoncharacteristicHyperplane.lean}{235}{Stafford38.NoncharacteristicHyperplane.canonical_isNoncharacteristic_annihilator},
\leandecl{Stafford38/NoncharacteristicHyperplane.lean}{331}{Stafford38.NoncharacteristicHyperplane.canonicalSupport_conormal_subset_zeroSection}}}


% Additional consequence declarations retained from Appendix B.
\AIadd{\formalref{\leandecl{Stafford38/LocalizationCorollaries.lean}{31}{Stafford38.LocalizationCorollaries.s38_rightOreLocalization}}}
\AIadd{\formalref{\leandecl{Stafford38/LeftHandedCorollary.lean}{23}{Stafford38.LeftHandedCorollary.leftHanded_of_universalStatement}}}
\AIadd{\formalref{\leandecl{Stafford38/LocalizedDifferentialCorollaries.lean}{83}{Stafford38.LocalizedDifferentialCorollaries.s38_unconditional_localized_differential}}}
\AIadd{\formalref{\leandecl{Stafford38/EvolutionaryCorollary.lean}{36}{Stafford38.Evolution.evolutionaryCorollary}}}
\fi

\bibliographystyle{alpha}
\bibliography{references}
\end{document}
